\documentclass[11pt]{article}
\usepackage[a4paper,margin=25mm]{geometry}
\usepackage{amsmath,amssymb,amsthm,booktabs,longtable}
\usepackage[colorlinks=true,urlcolor=blue]{hyperref}
\newtheorem{theorem}{Theorem}
\newtheorem{lemma}[theorem]{Lemma}
\newtheorem{corollary}[theorem]{Corollary}
\newcommand{\ind}{\mathbf 1}
\newcommand{\PP}{\mathbb P}
\newcommand{\Om}{\Omega}
\title{Independent audit of the full eleven-term quotient Chen weight}
\author{Independent mathematical audit}
\date{8 October 2026}
\begin{document}
\maketitle

This note proves the entire combinatorial interface from the actual
eleven quotient sums to the summed $D_{1,2}$ inequality. It does not
verify the subsequent prime-distribution, switching, or numerical
main-term estimates. No formal-kernel verification is claimed.
The manuscript inspected is \texttt{imbalanced-chen-transfer.tex},
SHA256
\begin{center}\small\ttfamily
3eabb2089dcdf3efcc030385f01a49630c36f4ea592bf763e3d86e08cc2c611e.
\end{center}
Its quotient definitions and direct weight proof occur at source lines
465--605. The proof below is independent of all old numerical PASS files.

\section{Definitions and exact parameter guards}
Let $N>1$ be real, and suppose
\begin{equation}\label{eq:guards}
 0<a<h<d<\frac13,\qquad h+d>\frac13,\qquad 4a+d<1.
\end{equation}
The manuscript has
\[
 a=\frac{100}{1149},\qquad h=\frac{50}{309},\qquad
 d=\frac{183}{766}=\frac12-3a.
\]
For this tuple the positive guard margins are
\[
 h-a=\frac{2950}{39449},\quad
 d-h=\frac{18247}{236694},\quad
 \frac13-d=\frac{217}{2298},\quad
 h+d-\frac13=\frac{15949}{236694},\quad
 1-4a-d=\frac{949}{2298}.
\]
The special identity $d=1/2-3a$ is not used below.

For a positive integer $n$ set $P^-(1)=+\infty$ and
\[
 Q_b(y;n)=\ind_{b\mid n}\ind_{P^-(n/b)\ge y}.
\]
The second factor is evaluated only when $b\mid n$.
All prefix letters in the following sums are distinct primes.
The lower endpoints are inclusive and upper endpoints are strict:
\begin{align*}
Y_1&=Q_1(N^a;n),&Y_2&=Q_1(N^h;n),\\
Y_3&=\sum_{N^a\le p<N^{1/3}}Q_p(N^a;n),&
Y_4&=\sum_{N^a\le p<N^d}Q_p(N^a;n),\\
Y_5&=\sum_{N^a\le p_1<p_2<N^h}Q_{p_1p_2}(N^a;n),&
Y_6&=\sum_{N^a\le p_1<N^h\le p_2<N^d}Q_{p_1p_2}(N^a;n),\\
Y_7&=\sum_{N^d\le p_1<p_2<(N/p_1)^{1/2}}Q_{p_1p_2}(p_2;n),\\
Y_8&=\sum_{N^a\le p_1<N^{1/3}\le p_2<(N/p_1)^{1/2}}
 Q_{p_1p_2}(p_2;n),\\
Y_9&=\sum_{N^h\le p_1<N^d\le p_2<(N/p_1)^{1/2}}
 Q_{p_1p_2}\bigl((N/(p_1p_2))^{1/2};n\bigr),\\
Y_{10}&=\sum_{N^a\le p_1<p_2<p_3<p_4<N^h}
 Q_{p_1p_2p_3p_4}(p_2;n),\\
Y_{11}&=\sum_{N^a\le p_1<p_2<p_3<N^h\le p_4<N^d}
 Q_{p_1p_2p_3p_4}(p_2;n).
\end{align*}
Thus $Y_{11}$ always denotes the full fourth-prime range. Put
\begin{equation}\label{eq:weight}
 W=3Y_1+Y_2-Y_3-Y_4+Y_5+Y_6-2Y_7-Y_8-Y_9-Y_{10}-Y_{11}.
\end{equation}
Throughout, $\Om$ counts prime factors with multiplicity and $\Om(1)=0$.

\begin{theorem}\label{thm:pointwise}
Under \eqref{eq:guards}, every squarefree integer $1\le n\le N$ satisfies
\[
                  W(n)\le4\ind_{\Om(n)\le2}.
\]
Neither $n\ge N^{1-a}$ nor $(n,N)=1$ is required.
\end{theorem}

\section{Quotient support and exact tuple identities}
\begin{lemma}\label{lem:rough}
If $n$ is not $N^a$-rough, all eleven terms vanish. On their nonzero
support, even without squarefreeness, $\Om(n)\le M:=\lfloor1/a\rfloor$.
\end{lemma}
\begin{proof}
In every term other than $Y_9$, all selected prefix primes and all
remaining factors are at least $N^a$, immediately from the displayed
ranges and quotient cutoff. For $Y_9$,
\[
 p_1p_2<(Np_1)^{1/2}<N^{(1+d)/2},\qquad
 (N/(p_1p_2))^{1/2}>N^{(1-d)/4}>N^a.
\]
Its selected primes are also larger than $N^a$; a unit quotient causes
no exception. Thus any nonzero term implies roughness. If $n$ is rough,
$N^{a\Om(n)}\le n\le N$, which gives the claimed factor bound.
\end{proof}

Assume now that $n$ is squarefree and rough. Write its prime factors
$q_1<\cdots<q_k$, and put
\[
 r=\#\{i:q_i<N^h\},\quad s=\#\{i:q_i<N^d\},\quad
 t=\#\{i:q_i<N^{1/3}\}.
\]
Then $0\le r\le s\le t\le k\le M$.

\begin{lemma}\label{lem:identities}
The first six terms are exactly
\[
 (Y_1,Y_2,Y_3,Y_4,Y_5,Y_6)
   =\left(1,\ind_{r=0},t,s,\binom r2,r(s-r)\right).
\]
For $r\ge2$,
\begin{equation}\label{eq:fourth}
 Y_{10}=\binom{r-2}{2},\qquad Y_{11}=(r-2)(s-r),
\end{equation}
where $\binom j2=0$ for $j=0,1$. For $k\ge3$,
\begin{equation}\label{eq:pair}
 t\ge1,\qquad Y_7=\ind_{s=0},\qquad Y_8=\ind_{t=1}.
\end{equation}
\end{lemma}
\begin{proof}
After removing any selected distinct factors, the remaining quotient
is still $N^a$-rough. This proves the first-six identities by counting
single factors and pairs in the indicated disjoint intervals.

Consider a tuple tested at its second selected factor $p_2$. All factors
of $n$ below $p_2$ must be removed. There is only one selected factor
below $p_2$, namely $p_1$. Hence any admissible such tuple has
$p_1=q_1,p_2=q_2$. Conversely, removing $q_1,q_2$ and any additionally
selected factors leaves either a unit or factors larger than $q_2$.
For $Y_{10}$ this leaves two arbitrary factors among the remaining
$r-2$ low factors. For $Y_{11}$ it leaves one such low factor and one
of the $s-r$ factors in $[N^h,N^d)$. This proves \eqref{eq:fourth},
including the zero cases; it uses the full upper endpoint $N^d$.

For $k\ge3$, squarefreeness gives
\[
                 q_1q_2^2<q_1q_2q_3\le n\le N.
\]
Consequently the pair $(q_1,q_2)$ always satisfies the strict switching
upper endpoint. The same preceding argument makes it the only possible
pair in $Y_7$ or $Y_8$. Their range conditions are respectively
$q_1\ge N^d$ and $q_1<N^{1/3}\le q_2$, giving the two indicators.
If $t=0$, the first three distinct factors are all at least $N^{1/3}$,
and their product is strictly greater than $N$. Thus $t\ge1$.
\end{proof}

The argument forcing the first two factors does \emph{not} apply to
$Y_9$, whose quotient cutoff is different. No uniqueness assertion for
$Y_9$ is used anywhere in this proof.

\begin{lemma}\label{lem:critical}
If $k\ge3$, $r=0$, $s=1$, and $t\le2$, then $k=3$ and $Y_9\ge1$.
\end{lemma}
\begin{proof}
The count conditions give
\[
 q_1\ge N^h,\quad q_1<N^d\le q_2,\quad q_3\ge N^{1/3}.
\]
If $k\ge4$, also $q_4\ge N^{1/3}$, and
$n\ge q_1q_2q_3q_4\ge N^{h+d+2/3}>N$, a contradiction.
Therefore $k=3$. The pair $(q_1,q_2)$ satisfies its strict upper endpoint
by Lemma~\ref{lem:identities}. Its quotient is $q_3$, and
\begin{equation}\label{eq:critical-cutoff}
 q_1q_2q_3^2\ge N^{h+d+2/3}>N.
\end{equation}
Thus $q_3>(N/(q_1q_2))^{1/2}$, so this pair is counted in $Y_9$.
No lower bound on $n$ was used.
\end{proof}

\section{Exhaustive analytic case proof}
\begin{proof}[Proof of Theorem~\ref{thm:pointwise}]
Nonrough $n$ are settled by Lemma~\ref{lem:rough}. If $n=1$, then
$Y_1=Y_2=1$, all other terms vanish, and $W=4$.
For rough squarefree $n>1$ use the preceding notation.

First suppose $r\ge2$. Dropping only nonpositive switching terms and
using \eqref{eq:fourth} yields
\begin{align*}
 W&\le3-t-s+\binom r2+r(s-r)
              -\binom{r-2}{2}-(r-2)(s-r)\\
  &=s-t\le0.
\end{align*}
The binomial identity is valid also at $r=2,3$ under the stated convention.

If $k\le2$, the first-six upper bound suffices. If $Y_5+Y_6=0$, then
$W\le3Y_1+Y_2\le4$. Otherwise there are exactly two factors,
$Y_5+Y_6=1$, $Y_2=0$, and $Y_3=Y_4=2$, so $W\le0$.
This also settles any overlap with the preceding $r\ge2$ case.

It remains that $k\ge3$ and $r\le1$, with $t\ge1$.
All possible count cases are the following:
\begin{center}
\begin{tabular}{lll}
\toprule
Count case & Justified upper bound for $W$ & Reason it is nonpositive\\
\midrule
$r=1$ & $2-t-\ind_{t=1}$ & $t\ge1$\\
$r=0,\ s\ge2$ & $4-s-t$ & $t\ge s\ge2$\\
$r=s=0$ & $2-t-\ind_{t=1}$ & $Y_7=1$, $t\ge1$\\
$r=0,\ s=1,\ t\ge3$ & $3-t$ & $t\ge3$\\
$r=0,\ s=1,\ t\le2$ & $3-t-1-\ind_{t=1}$ & $t=1$ or $2$\\
\bottomrule
\end{tabular}
\end{center}
The bounds use the first-six identities, $Y_7,Y_8$ from
\eqref{eq:pair}, and $Y_9\ge1$ only in the final row, where
Lemma~\ref{lem:critical} supplies it. All other omitted terms are
nonpositive. The rows exhaust $r\in\{0,1\}$ and $s\ge r$.
This proves the theorem for every integer $k$ permitted by roughness.
\end{proof}

There is no unexamined continuous-exponent branch. Indeed, putting
$x_i=\log(q_i)/\log N$ maps any actual factorization into
\[
 a\le x_1<\cdots<x_k,\qquad \sum_i x_i\le1.
\]
The half-open categories $[a,h)$, $[h,d)$, $[d,1/3)$, $[1/3,\infty)$
have sizes $r,s-r,t-s,k-t$. Every point in each such region is covered
by the corresponding row above, since the proof uses only its category
inequalities, strict ordering, and the sum constraint. In the critical
row the exact linear bounds are
\[
 x_1+x_2+x_3+x_4\ge h+d+\tfrac23>1\quad(k\ge4),\qquad
 x_1+x_2+2x_3\ge h+d+\tfrac23>1\quad(k=3).
\]
These establish infeasibility and the $Y_9$ cutoff, respectively.
No rounding, mesh density, or conjectured continuity is involved.

For the fixed tuple $M=11$, there are
$\sum_{k=0}^{11}\binom{k+3}{3}=\binom{15}{4}=1365$ count modes.
The companion certificate lists them all. It discards $t=0,k\ge3$
by strict ordering, or a mode with category lower sum
\[
 L=r a+(s-r)h+(t-s)d+(k-t)/3>1
\]
by the sum constraint, and certifies the displayed upper bound on every
remaining mode. This finite arithmetic check corroborates the analytic
proof; the analytic lemmas, not sampled exponents, justify its rules.
Modes with $L\le1$ need not actually be realizable: retaining extra modes
is harmless because the same upper bound holds on the larger class.

\section{Square factors, endpoints, and the summed inequality}
Squarefreeness cannot be dropped from the pointwise statement. For
$N=128$, $n=125=5^3$, and the fixed tuple, the eleven values are
\[
 (Y_1,\ldots,Y_{11})=(1,1,1,0,0,0,0,0,0,0,0).
\]
Thus $W=3$ although $\Om(n)=3$. This is also an actual shifted-prime
example $N-n=3$, with $(n,N)=1$, in the manuscript's original bulk range.
It is exactly a large-prime-square exception, not a failure of the
squarefree theorem.

On rough support all tuple counts are bounded even with multiplicity.
A convenient absolute bound is
\begin{equation}\label{eq:uniform-weight}
 |W|\le C_M:=4+2M+6\binom M2+2\binom M4.
\end{equation}
There are at most $M$ distinct factors, and a selected prefix contains
distinct primes, so the right side bounds each family by all possible
subsets and adds its coefficient magnitude. For $M=11$, $C_M=1016$.

Let $\mathcal E_N$ be the rough squareful positive integers at most $N$.
Writing $z=N^a>1$, the elementary union bound gives
\begin{equation}\label{eq:square-error}
 |\mathcal E_N|\le\sum_{q\ge z,\ q\in\PP}\left\lfloor\frac N{q^2}\right\rfloor
 \le N\sum_{m\ge\lceil z\rceil}\frac1{m^2}
 \le\frac N{\lceil z\rceil-1}
 \le2N^{1-a}\quad(z\ge2).
\end{equation}
The integral bound in the penultimate step is valid since
$\lceil z\rceil\ge2$. Nonrough squareful integers have all terms zero.
Consequently, for every positive $n\le N$,
\[
 W(n)\le4\ind_{\Om(n)\le2}+C_M\ind_{n\in\mathcal E_N}.
\]

\begin{corollary}\label{cor:summed}
For even integers $N\ge4$, let
\[
 D_{1,2}(N)=\#\{P<N:P\in\PP,\ \Om(N-P)\le2\},\qquad
 \Upsilon_i=\sum_{\substack{P<N,\ P\in\PP\\(P,N)=1}}Y_i(N-P).
\]
Then, for the fixed tuple and $N^a\ge2$,
\begin{align*}
 4D_{1,2}(N)\ge{}&3\Upsilon_1+\Upsilon_2-\Upsilon_3-\Upsilon_4
                  +\Upsilon_5+\Upsilon_6\\
 &-2\Upsilon_7-\Upsilon_8-\Upsilon_9-\Upsilon_{10}-\Upsilon_{11}
                  -2032N^{1-a}.
\end{align*}
In particular the manuscript's $O(N^{1-a})$ version follows.
\end{corollary}
\begin{proof}
Sum the preceding pointwise bound over the restricted prime set. The
map $P\mapsto N-P$ is injective, so the exceptional contribution is at
most $C_M|\mathcal E_N|$. The restricted almost-prime count is at most
$D_{1,2}(N)$. Apply \eqref{eq:square-error}. For even $N\ge4$, $N$ is
not prime, so the manuscript's $P\le N$ convention gives the identical
index set. No term has $n=0$.
\end{proof}

The coprimality restriction is useful for later analytic interfaces but
unnecessary for this combinatorial argument. Keeping it as in the
manuscript causes no extra error here. The unit $n=1$ is included
correctly with weight $4$ and $\Om(1)=0$; the later transfer to prime
or prime-times-prime representations must remove at most this one
unit. The manuscript explicitly does so in its $D_b\ge D_{1,2}-B_b-1$
comparison. Prime squares have $\Om=2$ but are harmlessly included in
the squareful exceptional upper bound.

All equalities at $N^a,N^h,N^d,N^{1/3}$ are handled by the half-open
categories. The switching upper endpoint is strict because $q_2<q_3$.
The quotient cutoff is inclusive; the critical case actually clears it
strictly by \eqref{eq:critical-cutoff}. Thus there is no endpoint term
missing from this pointwise proof.

\section{Negative controls and exact scope of the audit}
The full $Y_{11}$ range is essential to the cancellation proved here.
For the exponent pattern $(.09,.10,.15,.23,.42)$ at the fixed tuple,
the full term gives $W=0$, whereas imposing the additional restriction
$p_3p_4<N^{1/2-2a}$ gives $W=1$. There is an actual integer realization:
\[
\begin{split}
(q_1,q_2,q_3,q_4,q_5)&=(521,1031,32771,8388617,4398046511119),\\
n&=649435861824529986468405575083,\\
N=n+3&=649435861824529986468405575086.
\end{split}
\]
Trial division verifies all five primes. Exact integer-power comparisons
verify the prescribed ranges, $(n,N)=1$, and $n\ge N^{1-a}$.
The full vector is $(1,0,4,4,3,3,0,0,0,0,1)$; truncation changes only
its last entry to zero. This refutes that deletion for the present
quotient weight. It does not refute Wu's original aggregate theorem
with its original sieve objects and correction argument.

Similarly the exact exponent pattern $(1/5,39/100,79/200)$ has $Y_9=2$.
Its full vector is $(1,1,1,1,0,0,0,1,2,0,0)$ and $W=-1$.
This prevents mistakenly replacing the actual $Y_9$ by a unique tuple;
the proof only uses nonnegativity or the required lower bound $1$.

The independent certificate is
\texttt{independent-global-chen-weight-audit.json}, generated by
\texttt{independent-global-chen-weight-audit.py}. It binds the source
hashes and enumerates all integer count modes with exact rational
guards. Earlier finite boundary samples are diagnostics only.
The present result closes the stated combinatorial and elementary
exception-counting interface. It supplies no estimate for the
$\Upsilon_i$ in primes, no distribution theorem for a moving residue,
no proof of the claimed global Chen coefficient, and no formalization
of the whole manuscript.
\end{document}
